\documentclass[aspectratio=169,11pt,t]{beamer}
% ============================================================
%  ksa-cs-v2: ksa-cs-lecture-template.tex 스타일 + 한글, v2 덱(16:9)용.
%    - 짙은 네이비 헤드라인 (섹션 이름 + 프레임마다 점)
%    - 꽉 찬 파란 제목 막대, 둥근 그림자 블록
%    - 한글: Noto Sans CJK KR (fontspec)
%  사용:
%    \documentclass[aspectratio=169,11pt,t]{beamer}
%    \input{<..>/theme/ksa-cs-v2.tex}
%  ksa-theme.tex 의 매크로(\kb \kq, ksablue/ksared 등)도 같은 이름으로 둔다
%  -> v2 덱 본문을 거의 그대로 옮길 수 있다.
% ============================================================

\usepackage{fontspec}
\usepackage{amsmath,amssymb}
\usepackage{tikz}
\usepackage{listings}
\usepackage{graphicx}
\usepackage{booktabs}
\usepackage{varwidth}
% enumitem 은 넣지 않는다 (beamer 글머리표가 사라진다)

\defaultfontfeatures{Ligatures=TeX}
\setmainfont{Noto Sans CJK KR}
\setsansfont{Noto Sans CJK KR}
\setmonofont{Noto Sans Mono CJK KR}[Scale=0.92]
\usefonttheme{professionalfonts}

% ---------- 기본 틀 ----------
\usetheme{default}
\useinnertheme[shadow=true]{rounded}
\setbeamertemplate{navigation symbols}{}
\setbeamersize{text margin left=1em,text margin right=1em}

% ---------- 색 (템플릿 원본 값) ----------
\definecolor{navydark}{HTML}{191959}
\definecolor{navydim}{HTML}{8888B0}
\definecolor{navytitle}{HTML}{262686}
\definecolor{navyblue}{HTML}{3333B2}
\definecolor{blockbodybg}{HTML}{E9E9F3}
\definecolor{examplegreen}{HTML}{006000}
\definecolor{examplebg}{HTML}{E5EFE5}
\definecolor{alertred}{HTML}{BF0000}
\definecolor{alertbg}{HTML}{F9E5E5}
% ksa-theme.tex 와 같은 이름 (본문 호환)
\colorlet{ksablue}{navytitle}
\colorlet{ksabluedk}{navydark}
\definecolor{ksagold}{HTML}{D6CBB1}
\definecolor{ksagolddk}{HTML}{A9976B}
\colorlet{ksared}{alertred}
\definecolor{ksagray}{HTML}{5B5B5B}

\setbeamercolor{structure}{fg=navyblue}
\setbeamercolor{alerted text}{fg=alertred}
\setbeamercolor{frametitle}{bg=navytitle,fg=white}
\setbeamercolor{title}{fg=white,bg=navytitle}
\setbeamercolor{block title}{bg=navytitle,fg=white}
\setbeamercolor{block body}{bg=blockbodybg,fg=black}
\setbeamercolor{block title example}{bg=examplegreen,fg=white}
\setbeamercolor{block body example}{bg=examplebg,fg=black}
\setbeamercolor{block title alerted}{bg=alertred,fg=white}
\setbeamercolor{block body alerted}{bg=alertbg,fg=black}
\setbeamercolor{item}{fg=navyblue}
\setbeamertemplate{itemize item}[circle]
\setbeamertemplate{itemize subitem}[triangle]

% ---------- 헤드라인: miniframes (섹션 이름 + 점) ----------
\useoutertheme[subsection=false]{miniframes}
\setbeamerfont{section in head/foot}{size=\tiny}
\setbeamercolor{section in head/foot}{bg=navydark,fg=white}
\setbeamercolor{mini frame}{fg=white}
\setbeamertemplate{mini frames}[circle]

% ---------- 제목 막대: 폭 꽉 채움, 두 줄까지 ----------
\setbeamerfont{frametitle}{size=\normalsize,series=\bfseries}
\setbeamertemplate{frametitle}{%
  \nointerlineskip
  \begin{beamercolorbox}[wd=\paperwidth,ht=2.7ex,dp=1.1ex,leftskip=.3cm,rightskip=.3cm]{frametitle}%
    \usebeamerfont{frametitle}\insertframetitle
  \end{beamercolorbox}}

% ---------- 쪽 번호: 오른쪽 아래 ----------
\setbeamertemplate{footline}{%
  \hfill
  {\color{navydim}\scriptsize\insertframenumber\,/\,\inserttotalframenumber}%
  \hspace{3mm}\vspace{2mm}%
}

% ---------- 표지: 템플릿처럼 헤드라인은 그대로, 가운데 둥근 그림자 상자에 제목 한 줄 ----------
% 덱에서:  {\setbeamertemplate{footline}{}\begin{frame}[c,noframenumbering]\titlepage\end{frame}}
\setbeamercolor{title}{fg=white,bg=navyblue}
\makeatletter
\setbeamertemplate{title page}{%
  \vfill
  \centering
  \begin{minipage}{0.84\paperwidth}
  \begin{beamercolorbox}[wd=\linewidth,sep=7pt,center,rounded=true,shadow=true]{title}
    {\usebeamerfont{title}\inserttitle\par}%
    \ifx\insertsubtitle\@empty\else\vskip0.4em{\usebeamerfont{subtitle}\insertsubtitle\par}\fi
  \end{beamercolorbox}
  \end{minipage}\par\vskip1.6em
  {\usebeamerfont{author}\color{navytitle}\insertauthor\par}\ifx\beamer@date\@empty\vskip0.3em\else\vskip1.3em\fi
  {\usebeamerfont{institute}\color{ksagray}\insertinstitute\par}%
  \ifx\beamer@date\@empty\else\vskip0.3em{\usebeamerfont{date}\color{ksagray}\insertdate\par}\fi
  \vfill
}
\makeatother
\setbeamerfont{title}{series=\bfseries,size=\Large}
\setbeamerfont{subtitle}{size=\normalsize}
\setbeamerfont{author}{size=\normalsize}
\setbeamerfont{institute}{size=\small}
\setbeamerfont{date}{size=\small}

% ksa-theme 호환: 덱이 \renewcommand{\ksacredit}{..} 를 해도 되게 (꼬리말에는 안 쓴다)
\newcommand{\ksacredit}{}

% ---------- 강조 ----------
\newcommand{\hl}[1]{\textcolor{alertred}{#1}}
\newcommand{\hlb}[1]{\textcolor{navyblue}{#1}}
\newcommand{\kb}[1]{\textbf{\color{ksablue}#1}}
\newcommand{\kq}[1]{\textbf{\color{ksared}#1}}

% ---------- 파이썬 코드 (템플릿과 같은 색) ----------
\definecolor{codekw}{HTML}{0033B3}
\definecolor{codebuiltin}{HTML}{BF0000}
\definecolor{codestr}{HTML}{067D17}
\definecolor{codecomment}{HTML}{8C8C8C}
\definecolor{codenum}{HTML}{1750EB}
\lstdefinestyle{py}{
  language=Python, basicstyle=\ttfamily\scriptsize,
  keywordstyle=\color{codekw}\bfseries, stringstyle=\color{codestr},
  commentstyle=\color{codecomment}\itshape, numberstyle=\color{codenum},
  emph={print,len,range,sum,input,int,str,float,list,dict,append,remove,sorted,abs,min,max},
  emphstyle=\color{codebuiltin}, showstringspaces=false, keepspaces=true,
  columns=fixed, basewidth=0.52em, aboveskip=1mm, belowskip=1mm, xleftmargin=2mm,
  breaklines=true
}
\lstset{style=py}

\usepackage{fontawesome5}
\newcommand{\paperlink}[1]{\href{#1}{\textcolor{ksablue}{\faLink}}}
% GLM 보충 덱 (2026-11-27): 수업 덱 ../1_glm 에서 뺀 지수족 · 세 가정 · 유도 · a'=평균 · a''=볼록.
% 숫자는 data.py (assert), 그림은 make_figs.py.
\setbeamerfont{author}{size=\normalsize}
\setbeamerfont{institute}{size=\small}
\graphicspath{{figs/}{../../../assets/}}
\usetikzlibrary{positioning,arrows.meta,calc}

\renewcommand{\ksacredit}{Machine Learning and Deep Learning \,\textperiodcentered\, GLM 보충}
\title{GLM 보충: 지수족}
\subtitle{기울기가 모두 $\hat y - y$ 인 이유 (수업 밖, 궁금한 사람만)}
\author{Machine Learning and Deep Learning}
\institute{KSA 2027-Spring}
\date{Sumin Han (suminhan@ksa.hs.kr)}
\begin{document}
{\setbeamertemplate{footline}{}\begin{frame}[c,noframenumbering]\titlepage\end{frame}}

\section{질문}
% --- 우연일까
\begin{frame}{네 모델, 기울기가 모두 같다: 우연일까?}
\vfill
\centering
\renewcommand{\arraystretch}{1.5}
{\large\begin{tabular}{l|c|c|c}
모델 & 예측 $\hat y$ & 손실 & $\partial\,\text{loss} / \partial z$\\ \hline
선형회귀 & $z$ & $\tfrac12 (y - \hat y)^2$ & \kq{$\hat y - y$}\\
로지스틱회귀 & $\sigma(z)$ & $-\big[y\log\hat y + (1-y)\log(1-\hat y)\big]$ & \kq{$\hat y - y$}\\
포아송 회귀 & $e^{z}$ & $\hat y - y\log\hat y$ & \kq{$\hat y - y$}\\
소프트맥스 회귀 & $\mathrm{softmax}(z)$ & $-\log\hat y_{\text{정답}}$ & \kq{$\hat y - y$}
\end{tabular}}\\[0.8em]
{\large 함수도 손실도 다 다른데 기울기는 \kb{모두} $\hat y - y$ \;$\to$\; 네 모델을 묶는 \kb{일반화된 식}이 있다}\\[0.3em]
{\small 이 보충 덱: 네 분포가 같은 꼴 (지수족) $\to$ 기울기가 왜 같은지, 손실이 왜 볼록인지}
\vfill
\note{수업에서 본 네 모델을 다시 나란히 놓아 봅시다. 예측 함수는 z 그대로, 시그모이드, e의 z, 소프트맥스로 다 다릅니다. 손실도 오차 제곱, 교차 엔트로피, 포아송 손실, 다중 교차 엔트로피로 다 다릅니다. 그런데 손실을 z로 미분하면 넷 다 와이햇 빼기 y입니다. 우연일까요? 아닙니다. 이 네 모델은 한 가지 식에서 나옵니다. 그 일반화된 식이 일반화 선형 모델, GLM입니다. 이름이 뜻을 말해 줍니다. 선형회귀는 y가 직선 근처에 정규분포로 흩어져 있다고 봅니다. 직선, Linear는 그대로 두고, 정규분포가 아닌 y까지 다룰 수 있게 넓힌 것이 일반화, Generalized입니다.}
\end{frame}

\section{지수족}
% --- 지수족
\begin{frame}{네 분포가 같은 꼴로 써진다: 지수족}
\vfill
\centering
{\Large\[ p(y;z) = b(y)\,\exp\!\big(\,\underbrace{z}_{\text{natural parameter}}\, \underbrace{T(y)}_{\text{sufficient statistic}} \;-\; \underbrace{a(z)}_{\text{log partition}}\,\big) \]}\vskip0.1em
\renewcommand{\arraystretch}{1.12}
{\large\begin{tabular}{c|l}
$p(y;z)$ & 손잡이를 $z$ 로 맞춘 분포에서 \kb{$y$ 가 나올 확률}\\
$y$ & 실제로 나온 값 (골 수, 승패, 점수)\\
$z$ & 분포를 정하는 \kb{손잡이} (자연 매개변수)\\
$T(y)$ & $y$ 에서 뽑아 쓰는 값 (오늘은 $y$ 그대로)\\
$b(y)$ & $y$ 만 들어 있는 부분\\
$a(z)$ & $z$ 만 들어 있는 부분 $\to$ 확률의 합을 1 로 맞춘다
\end{tabular}}\\[0.3em]
{\large 정규 · 베르누이 · 포아송 · 범주형이 모두 이 꼴 $=$ \kb{지수족}}
\vfill
\note{네 모델의 분포, 정규, 베르누이, 포아송, 범주형은 정리하면 모두 이 꼴로 쓸 수 있습니다. b of y 곱하기 exp z 곱하기 T of y 빼기 a of z. 기호를 하나씩 봅시다. p of y 세미콜론 z는 손잡이를 z로 맞춘 분포에서 y가 나올 확률입니다. 세미콜론 뒤는 설정값이라는 뜻입니다. y는 실제로 나온 값, 골 수나 승패나 점수입니다. z는 분포를 정하는 손잡이로, 돌리면 평균이 바뀝니다. b of y는 y만, a of z는 z만 들어 있는 부분입니다. 이 꼴로 쓸 수 있는 분포들을 지수족이라고 부릅니다. z를 자연 매개변수, T of y를 충분 통계량, a of z를 로그 분배 함수라고 부릅니다. 이름은 어렵지만 역할은 간단합니다. T, a, b 세 가지를 정하면 분포의 가족 하나가 정해지고, 그 안에서 z를 바꾸면 평균이 다른 분포들이 나옵니다. 예를 들어 베르누이 가족에서 z를 바꾸면 앞면 확률이 다른 동전들이 나옵니다. e의 마이너스 a of z는 확률을 다 더했을 때 1이 되도록 맞추는 몫입니다. 오늘 보는 분포는 범주형만 빼고 T of y가 y 그대로입니다. 이제 실제 분포를 이 꼴에 넣어 봅시다.}
\end{frame}

% --- 지수족
\begin{frame}{분포들을 이 꼴에 넣으면}
\vfill
\centering
{\Large\[ p(y) = \textcolor{gray}{b(y)}\,\cdot\,\exp\!\big(\,\textcolor{ksared}{z}\cdot y \;-\; \textcolor{ksablue}{a(z)}\,\big) \]}\vskip0.3em
\renewcommand{\arraystretch}{1.4}
{\large\begin{tabular}{l|c|c|c|c}
분포 & 원래 식 & \textcolor{gray}{$b(y)$}: $y$ 만 & \textcolor{ksared}{$z$}: $y$ 에 곱함 & \textcolor{ksablue}{$a(z)$}: 빼는 몫\\ \hline
포아송 & $\dfrac{\lambda^{y} e^{-\lambda}}{y!}$ & $\dfrac{1}{y!}$ & \textcolor{ksared}{$\log\lambda$} & \textcolor{ksablue}{$\lambda$}\\
베르누이 & $\phi^{\,y}(1-\phi)^{1-y}$ & $1$ & \textcolor{ksared}{$\log\frac{\phi}{1-\phi}$} & \textcolor{ksablue}{$-\log(1-\phi)$}\\
정규 ($\sigma{=}1$) & $\frac{1}{\sqrt{2\pi}}e^{-(y-\mu)^2/2}$ & $\frac{1}{\sqrt{2\pi}}e^{-y^2/2}$ & \textcolor{ksared}{$\mu$} & \textcolor{ksablue}{$\tfrac12\mu^2$}
\end{tabular}}\\[0.5em]
{\large 포아송 확인: $\lambda^{y} = e^{\,y\log\lambda}$ 이라서 $\frac{\lambda^{y} e^{-\lambda}}{y!} = \frac{1}{y!}\exp(\textcolor{ksared}{\log\lambda}\cdot y - \textcolor{ksablue}{\lambda})$ \quad 이 꼴 $=$ \kb{지수족}}\\[0.2em]
{\small $T(y) = y$ 라서 $z^\top T(y) = z \cdot y$ \;·\; 범주형은 $z$ 가 벡터 (뒤에서)}
\vfill
\note{앞에서 본 분포 식들은 모양이 제각각입니다. 그런데 정리하면 모두 방금 본 꼴로 쓸 수 있습니다. T of y는 y 그대로라서 z 곱하기 y로 적었습니다. 회색 b of y는 y만 들어 있는 부분, 빨간 z는 y에 곱해지는 부분, 파란 a of z는 빼 주는 부분입니다. 가장 쉬운 포아송으로 확인해 봅시다. 람다의 y제곱은 e의 y log 람다와 같습니다. 그러면 람다의 y제곱 곱하기 e의 마이너스 람다 나누기 y 팩토리얼은, 1 나누기 y 팩토리얼 곱하기 exp log 람다 곱하기 y 빼기 람다입니다. b of y는 1 나누기 y 팩토리얼, z는 log 람다, a는 람다입니다. 베르누이는 y log 파이 더하기 1 빼기 y log 1 빼기 파이를 y로 묶으면 z가 로그 오즈, 정규분포는 제곱을 전개하면 z가 뮤 그대로입니다. 앞 쪽의 붓꽃, 범주형 분포도 지수족인데 z가 숫자 하나가 아니라 클래스마다 하나씩인 벡터라서 표에서는 뺐고, 뒤의 소프트맥스에서 다시 나옵니다. 포아송이면 z가 평균 람다를 정하고, 베르누이면 앞면 확률 파이를 정합니다.}
\end{frame}

% --- GLM 세 가정
\begin{frame}{GLM: 세 가지만 정하면 된다}
\vfill
\centering
\renewcommand{\arraystretch}{1.6}
{\Large\begin{tabular}{rl}
\textbf{1.} & $y \mid x \;\sim\;$ \kb{지수족}$(z)$ \quad {\large ($y$ 의 분포를 고른다)}\\
\textbf{2.} & 예측 $\hat y = E[\,y \mid x\,]$ \quad {\large (그 분포의 \kb{평균})}\\
\textbf{3.} & $z = w_0 + w_1 x$ \quad {\large (매개변수 $=$ \kb{직선}, 가정이라기보다 \kq{설계 선택})}
\end{tabular}}\\[0.8em]
{\large 2번 확인 (로지스틱): \;$E[\,y \mid x\,] = 0 \cdot P(y{=}0) + 1 \cdot P(y{=}1) = P(y{=}1 \mid x) = \hat y$ \;\checkmark}\\[0.3em]
{\small GLM = Generalized Linear Model (일반화 선형 모델) \quad $x$ 가 여럿이면 $z = w \cdot x + b$}
\vfill
\note{GLM은 모델을 만드는 요리법이고, 세 가지를 정합니다. 첫째, x가 주어졌을 때 y가 지수족의 어떤 분포를 따른다고 고릅니다. 자연 매개변수는 z입니다. 둘째, 예측 와이햇은 그 분포의 평균, x가 주어졌을 때 y의 기댓값으로 합니다. 로지스틱회귀로 확인해 봅시다. y가 0 아니면 1이니 기댓값은 0 곱하기 0일 확률 더하기 1 곱하기 1일 확률, 곧 1일 확률입니다. 로지스틱회귀의 와이햇이 바로 이 확률이었으니 조건을 맞춥니다. 셋째, 자연 매개변수 z를 x의 직선으로 둡니다. 강의 노트는 이것이 가정이라기보다 설계 선택이라고 말합니다. 왜 꼭 직선이어야 하는지 근거가 있는 것은 아니지만, 이렇게 두면 학습이 아주 쉬워집니다. 이 세 가지로 지금까지의 모델이 어떻게 나오는지 봅시다.}
\end{frame}

% --- 유도: 선형회귀 · 로지스틱회귀
\begin{frame}{GLM 으로 유도하면: 선형회귀와 로지스틱회귀 (포아송도 같다)}
\vskip-0.3em
\begin{columns}[T]
\begin{column}{0.48\textwidth}
\centering{\large\kb{$y$ 가 실수: 정규}}
{\large\begin{align*}
\hat y &= E[\,y \mid x\,] && \text{2번}\\
&= \mu && \text{정규의 평균}\\
&= z && z = \mu\\
&= w_0 + w_1 x && \text{3번}
\end{align*}}
\centering{\large $=$ \kb{선형회귀}}
\end{column}
\begin{column}{0.5\textwidth}
\centering{\large\kb{$y$ 가 0 / 1: 베르누이}}
{\large\begin{align*}
\hat y &= E[\,y \mid x\,] && \text{2번}\\
&= \phi && \text{베르누이의 평균}\\
&= \dfrac{1}{1 + e^{-z}} && z = \text{로그 오즈를 풂}\\
&= \dfrac{1}{1 + e^{-(w_0 + w_1 x)}} && \text{3번}
\end{align*}}
\centering{\large $=$ \kb{로지스틱회귀}}
\end{column}
\end{columns}\vskip0.1em
\centering
{\large 시그모이드는 S 자라서 \textbf{고른} 것이 아니라, $y$ 가 0 / 1 이라서 \kb{나온} 것}\\[0.3em]
{\small 체스: 백이 200점 높으면 $z = 0.09 + 0.0047 \times 200 = 1.03$ \;$\to$\; $\sigma(1.03) = 74\%$}\\[0.1em]
{\small $z$ 를 평균으로 바꾸는 $g$ $=$ canonical response function (정규: 그대로, 베르누이: $\sigma$)}
\vfill
\note{세 가지 재료로 선형회귀를 유도해 봅시다. y가 실수라서 정규분포를 고릅니다. 와이햇은 2번에 따라 y의 기댓값, 정규분포의 기댓값은 뮤, 지수족 꼴에서 뮤는 z 그대로였고, 3번에 따라 z는 w0 더하기 w1 x입니다. 와이햇이 직선, 선형회귀가 나왔습니다. 로지스틱회귀도 같은 순서입니다. y가 0 아니면 1이라서 베르누이를 고릅니다. 와이햇은 y의 기댓값, 베르누이의 기댓값은 파이, 지수족 꼴에서 z 이퀄 로그 오즈를 풀면 파이는 1 나누기 1 더하기 e의 마이너스 z, 3번에 따라 z는 직선입니다. 로지스틱회귀가 나왔습니다. 지난 시간에는 시그모이드를 0과 1 사이의 S 자 함수라서 골랐습니다. 강의 노트의 말대로, 시그모이드가 왜 나왔는지에 대한 답이 여기 있습니다. 고른 게 아니라, y가 0 아니면 1이라서 베르누이를 고른 순간 저절로 나온 함수입니다. 체스로 확인하면, 백이 200점 높을 때 z는 1.03이고 시그모이드를 씌우면 74퍼센트, 지난 시간의 숫자와 같습니다. z를 평균으로 바꾸는 이 함수 g를 정준 반응 함수, canonical response function이라고 부릅니다. 포아송도 같습니다. 와이햇은 평균 람다, z는 log 람다였으니 람다는 e의 z. 앞에서 음수를 피하려고 씌운 e의 z도 고른 게 아니라 포아송에서 나온 함수였습니다.}
\end{frame}

% --- 손실 = -log p
\begin{frame}{손실은 어디서 오나: $-\log p$}
\vfill
\centering
{\LARGE\[ \text{loss} = -\log p(y;z) = \underbrace{-\,z\,y + a(z)}_{\text{\normalsize 분포마다 이 두 항}} \;+\; \text{const} \]}\vskip0.4em
\renewcommand{\arraystretch}{1.45}
{\large\begin{tabular}{l|l}
정규 & $\tfrac12 (y - \hat y)^2$ \;=\; \kb{오차 제곱}\\
베르누이 & $-\big[\,y\log\hat y + (1-y)\log(1-\hat y)\,\big]$ \;=\; \kb{교차 엔트로피}\\
포아송 & $\hat y - y\log\hat y$
\end{tabular}}\\[0.6em]
{\small 최대 가능도: 데이터가 나올 확률을 가장 크게 $=$ $-\log$ 확률의 합을 가장 작게 \quad 범주형 $\to$ $-\log\hat y_{\text{정답}}$}
\vfill
\note{손실은 어디서 왔을까요? GLM에서는 손실을 정하는 것도 요리법의 일부입니다. 학습 데이터가 나올 확률을 가장 크게 하는 w를 찾습니다. 이걸 최대 가능도라고 합니다. 확률을 크게 하는 것은 마이너스 로그 확률을 작게 하는 것과 같으니, 손실은 마이너스 log p입니다. 지수족 꼴에 넣으면 마이너스 z y 더하기 a of z, 그리고 w와 상관없는 상수입니다. 정규분포에서 계산하면 2분의 1 y 빼기 와이햇의 제곱, 1주차의 오차 제곱입니다. 베르누이에서 계산하면 지난 시간의 교차 엔트로피입니다. 그러니까 선형회귀가 오차 제곱을, 로지스틱회귀가 교차 엔트로피를 쓰는 건 따로 고른 게 아니라, 분포를 고르면 같이 정해지는 짝입니다. 포아송이면 와이햇 빼기 y log 와이햇입니다.}
\end{frame}

% --- 기울기 유도
\begin{frame}{기울기가 같은 이유: $a'(z)$ 가 곧 평균}
\vfill
\centering
{\large\begin{align*}
\frac{\partial\,\text{loss}}{\partial z} &= -y + a'(z) && \text{loss} = -z y + a(z) + \text{const}\\[0.3em]
&= -y + \textcolor{ksared}{\hat y} && a'(z) = \text{평균} = \hat y\\[0.3em]
\frac{\partial\,\text{loss}}{\partial w_1} &= (\hat y - y)\,x && z = w_0 + w_1 x
\end{align*}}\vskip0.1em
\renewcommand{\arraystretch}{1.2}
{\large\begin{tabular}{l|c|c|c}
 & $a(z)$ & $a'(z)$ & 평균\\ \hline
베르누이 & $\log(1 + e^{z})$ & $\frac{e^{z}}{1 + e^{z}} = \sigma(z)$ & $\phi$ \;\checkmark\\
정규 & $\tfrac12z^2$ & $z$ & $\mu$ \;\checkmark\\
포아송 & $e^{z}$ & $e^{z}$ & $\lambda$ \;\checkmark
\end{tabular}}
\vfill
\note{이제 네 모델의 기울기가 왜 모두 와이햇 빼기 y였는지 보입니다. 손실 마이너스 z y 더하기 a of z를 z로 미분하면 마이너스 y 더하기 a 프라임 z입니다. 여기서 놀라운 사실이 있습니다. a 프라임 z가 언제나 그 분포의 평균입니다. 아래 표로 확인해 봅시다. 베르누이의 a는 log 1 더하기 e의 z, 미분하면 시그모이드 z, 평균 파이입니다. 정규분포의 a는 2분의 1 z제곱, 미분하면 z, 평균 뮤입니다. 포아송의 a는 e의 z, 미분해도 e의 z, 평균 람다입니다. 평균은 곧 예측 와이햇이니, z로 미분한 값은 와이햇 빼기 y입니다. 마지막으로 z가 w0 더하기 w1 x이니 연쇄 법칙으로 x를 곱하면, w1으로 미분한 값은 와이햇 빼기 y 곱하기 x입니다.}
\end{frame}

% --- 한 번 더 미분: a'' = 분산 → 볼록
\begin{frame}{한 번 더 미분하면: $a''(z)$ 는 분산 $\to$ 손실은 볼록}
\vfill
\centering
{\Large\[ \frac{\partial^2\,\text{loss}}{\partial z^2} = a''(z) = \text{Var}(y) \;\ge\; 0 \]}\vskip0.4em
\renewcommand{\arraystretch}{1.2}
{\large\begin{tabular}{l|c|c}
 & $a''(z)$ & 분산\\ \hline
베르누이 & $\sigma(z)\,\big(1 - \sigma(z)\big)$ & $\phi(1-\phi)$ \;\checkmark\\
정규 & $1$ & $\sigma^2 = 1$ \;\checkmark\\
포아송 & $e^{z}$ & $\lambda$ \;\checkmark
\end{tabular}}\\[0.5em]
{\large 기울기가 늘 0 이상 \;$\to$\; 손실은 \kb{그릇 모양} (볼록) \;$\to$\; 경사하강법이 \kb{하나뿐인 바닥}으로 간다}\\[0.3em]
{\small $z$ 가 $w$ 의 직선이라 $w$ 에 대해서도 볼록 \quad (CS229 강의 4)}
\vfill
\note{한 번 더 미분해 봅시다. 손실을 z로 두 번 미분하면 a 두 번 미분이고, 이건 언제나 y의 분산입니다. 표로 확인하면 베르누이는 시그모이드 곱하기 1 빼기 시그모이드, 파이 곱하기 1 빼기 파이로 베르누이의 분산입니다. 정규는 1, 시그마를 1로 두었으니 분산 1입니다. 포아송은 e의 z, 람다이고 포아송은 분산도 람다입니다. 분산은 음수가 될 수 없으니 두 번 미분한 값이 언제나 0 이상, 손실은 그릇 모양, 볼록 함수입니다. z가 w의 직선이니 w에 대해서도 볼록합니다. 그래서 GLM은 어느 분포를 고르든 바닥이 하나뿐이고, 경사하강법이 엉뚱한 곳에 갇히지 않습니다. 1주차 선형회귀와 지난 시간 로지스틱회귀의 손실이 그릇 모양이었던 이유입니다. 앤드류 응의 강의 4강에서 다루는 지수족의 성질입니다.}
\end{frame}

% --- 한 표
\begin{frame}{정리: 분포만 바꾸면 모델이 나온다}
\vfill
\centering
\renewcommand{\arraystretch}{1.7}
{\large\begin{tabular}{l|l|c|c|l}
$y$ & 분포 & 매개변수 $z$ & 예측 $\hat y = g(z)$ & 모델\\ \hline
실수 & 정규 & $\mu$ & $z$ & 선형회귀\\
0 / 1 & 베르누이 & $\log\frac{\phi}{1-\phi}$ & $\sigma(z)$ & 로지스틱회귀\\
개수 & 포아송 & $\log\lambda$ & $e^{z}$ & 포아송 회귀\\
$K$ 개 중 하나 & 범주형 & (벡터) & $\mathrm{softmax}(z)$ & 소프트맥스 회귀
\end{tabular}}\\[1.0em]
{\large 직선 $z$ 는 그대로 \;·\; 바깥 함수 $g$ 만 바뀐다}\\[0.3em]
{\small $g$ = response function, \;$g^{-1}$ = \kb{link function} (로지스틱: 로그 오즈, 포아송: $\log$)}
\vfill
\note{지금까지를 한 표로 모읍니다. y가 실수면 정규분포, z는 뮤, 예측은 z 그대로, 선형회귀입니다. 0과 1이면 베르누이, z는 로그 오즈, 예측은 시그모이드 z, 로지스틱회귀입니다. 개수면 포아송, z는 log 람다, 예측은 e의 z, 포아송 회귀입니다. 클래스가 여럿이면 범주형, 예측은 소프트맥스, 소프트맥스 회귀입니다. 모든 줄에서 직선 z는 똑같고, 바깥에 씌우는 함수 g만 바뀝니다. 이 g를 반응 함수, 거꾸로 평균을 z로 바꾸는 함수를 링크 함수라고 부릅니다. 로지스틱회귀의 링크 함수는 로그 오즈, 포아송 회귀의 링크 함수는 로그입니다.}
\end{frame}

% --- 참고 문헌
\begin{frame}{참고 문헌}
\vfill
\small
\begin{itemize}
\item Ng, A., \& Ma, T. CS229 Lecture Notes, Chapter 3: Generalized Linear Models. Stanford. \paperlink{https://cs229.stanford.edu/main_notes.pdf}
\item Ng, A. (2018). Lecture 4: Perceptron \& Generalized Linear Model, CS229 Autumn 2018. \paperlink{https://www.youtube.com/watch?v=iZTeva0WSTQ}
\item McCullagh, P., \& Nelder, J. A. (1989). \textit{Generalized Linear Models} (2nd ed.). Chapman \& Hall.
\item Jürisoo, M. International football results from 1872 (국가대표 경기 결과, CC0). Elo 는 World Football Elo Ratings (eloratings.net) 방식으로 직접 계산. \paperlink{https://github.com/martj42/international_results}
\item Fisher, R. A. (1936). The use of multiple measurements in taxonomic problems. \textit{Annals of Eugenics}, 7(2).
\item lichess.org open database, 2013년 1월 (CC0). \paperlink{https://database.lichess.org}
\end{itemize}
\vfill
\note{오늘 내용은 스탠퍼드 CS229 강의 노트의 3부, 일반화 선형 모델을 바탕으로 했습니다. 앤드류 응의 2018년 강의 4강 영상에서 같은 내용을 볼 수 있습니다. 월드컵 데이터는 1872년부터의 국가대표 경기 결과 공개 데이터이고, Elo 점수는 그 경기 기록으로 직접 계산했습니다. 붓꽃은 피셔의 1936년 데이터, 체스는 lichess 공개 데이터베이스입니다.}
\end{frame}
\end{document}
